\documentclass[12pt, letterpaper]{article}
\usepackage{graphicx}
\usepackage[margin=1in]{geometry}
\usepackage{amsmath, amssymb, amsthm, mathrsfs}
\usepackage{fancyhdr}
\usepackage{tikz}
\usepackage{enumerate}
\usepackage{cancel}
\usetikzlibrary{positioning}
\usepackage[utf8]{inputenc}
\usepackage{xcolor}
\usepackage{framed}
\title{HW 1.3}
\author{Your Name}
\date{\today}
\pagestyle{fancy}
\renewcommand{\headrulewidth}{0pt}
\renewcommand{\footrulewidth}{0pt}
\fancyhf{}
\rhead{
    Zack Abdirashid\\
    4/21/26 \\
    MATH 405
}
\rfoot{Page \thepage}
\begin{document}
\paragraph{\underline{5.1} \\ \\} 
\textbf{\phantom{1}6.} \begin{tabular}[t]{l} 
    How many ways are there to pick a man and a woman who are not married to each \\
    other from a group of $n$ married couples?
\end{tabular} \\ \\
\indent \indent \begin{tabular}[t]{l}
    We can represent the married couples as the pairs $(h_1, w_1), (h_2, w_2), \dots, (h_n, w_n)$. \\ 
    For each husband $h_i$, we can choose any of the $n-1$ women who are not married \\
    to him. That is, 
    \end{tabular}
    \begin{align*}
        h_1 &: (h_1, w_2), (h_1, w_3), \dots, (h_1, w_n) \\
        h_2 &: (h_2, w_1), (h_2, w_3), \dots, (h_2, w_n) \\
        \vdots  \\
        h_n &: (h_n, w_1), (h_n, w_2), \dots, (h_n, w_{n-1}) \\
        \end{align*}
\indent \indent \begin{tabular}[t]{l}
    This is also analogous for the women, so we need not consider both cases to avoid\\
     overcounting. Since there are $n$ husbands, the total number of ways are \boxed{\text{$n(n-1)$}}
    \end{tabular} \\

\textbf{12.} \begin{tabular}[t]{l} 
    How many six-letter "words" (sequences of letters with repetition) are there in which \\
    the first and last letter are vowels? In which vowels appear only (if at all) as the first \\
    and last letter?
    \end{tabular} \\ \\
    \indent \indent \begin{tabular}[t]{l}
        (i) Consider the set of vowels $\{a, e, i, o, u\}$. Since the first and last letter can only be vo- \\
        wels, there are 5 choices for both. Due to repetition, the middle 4 letters have 26 choices \\
        each. So, there $5^{2} \cdot 26^{4} = \boxed{\text{11,424,400}}$ six-letter words for this case.
        \end{tabular}
        \indent \indent \begin{tabular}[t]{l}
            (ii) Now, words don't have to contain vowels, but they must be in the first and last \\
            position if they do. So, the choices for each letter are
            \end{tabular}
            \[
                \underset{\text{\small any letter}}{\overset{26}{\underline{\hspace{1cm}}}} \cdot
                \underset{\text{\small no vowels}}{\overset{21}{\underline{\hspace{1cm}}}} \cdot
                \underset{\text{\small no vowels}}{\overset{21}{\underline{\hspace{1cm}}}} \cdot
                \underset{\text{\small no vowels}}{\overset{21}{\underline{\hspace{1cm}}}} \cdot
                \underset{\text{\small no vowels}}{\overset{21}{\underline{\hspace{1cm}}}} \cdot
                \underset{\text{\small any letter}}{\overset{26}{\underline{\hspace{1cm}}}}
            \]
            \indent \indent \indent \begin{tabular}[t]{l}
                So, the total number of six-letter words are $26^{2} \cdot 21^{4} = \boxed{\text{131,469,156}}$
                \end{tabular} \\ \\
        \indent \indent \begin{tabular}[t]{l}
            Comment on these possible answers to the second part: \\
            \indent \indent (a) $5^{2}21^{4}$: Incorrect because it doesn't take into account that the word doesn't have \\
            \indent \indent \indent to contain vowels. So, the first and last letters have more than 5 choices. \\
            \indent \indent (b) $5^{2}26^{4}$: This is incorrect because of the reason for (a) along with the middle letters \\
            \indent \indent \indent being restricted to not being vowels. \\
            \indent \indent (c) $21^{2}26^{4}$: Incorrect because it flips the rules for the first/last letters and the middle \\
            \indent \indent \indent letters. \\
            \indent \indent (d) $26^{6} - 21^{2}26^{4}$: Incorrect because this difference essentially says \\
            \indent \indent \indent (all possible words) $-$ (words without vowels in the 1st and last position). \\
            \indent \indent \indent Due to that $26^{4}$ factor, this violates the rules for (ii) because the result can have \\
            \indent \indent \indent words with vowels in the middle, which is not allowed.
            \end{tabular} \\ \\
    \textbf{14.} \begin{tabular}[t]{l} 
        How many different numbers can be formed by various arrangements of the six digits \\
        $1, 1, 1, 1, 2, 3$?
        \end{tabular} \\ \\
        \indent \indent \begin{tabular}[t]{l}
            Consider the choices for the 2 and 3. Since there are 6 digits, there are 6 ways to place\\
            either of them. Then, the remaining of the two has 5 choices. Once we've place the 2\\
            and 3, there's only one way to place the 1's in the remaining spots since they're identical. \\
            So,
            \end{tabular} \\
            \[
            \underset{\text{\small ways to place choice of 2 or 3}}{\overset{6}{\underline{\hspace{1cm}}}} \cdot
            \underset{\text{\small ways to place remaining choice of 2 or 3}}{\overset{5}{\underline{\hspace{1cm}}}}
            \]

            \indent \indent \begin{tabular}[t]{l}
                There are $6 \cdot 5 = \boxed{\text{30}}$ different numbers that can be formed.
                \end{tabular} \\
    \textbf{22.} \begin{tabular}[t]{l} 
        How many ternary (0, 1, 2) sequences of length 10 are there without any pair of consecutive \\
        digits the same?
    \end{tabular} \\ \\
    \begin{tabular}[t]{l}
        The initial number will have 3 choices, but the number after it will have 2. The number after that \\
       will also have 2 choices. Thus, every subsequent number after the first will have 2 choices to ensure \\ 
       no consecutive digits are the same. So,
        \end{tabular}
        \[
            \underset{\text{\small 1st \#}}{\overset{3}{\underline{\hspace{1cm}}}} \cdot
            \underset{\text{\small 2nd \#}}{\overset{2}{\underline{\hspace{1cm}}}} \cdot
            \underset{\text{\small 3rd \#}}{\overset{2}{\underline{\hspace{1cm}}}} \cdot
            \underset{\text{\small 4th \#}}{\overset{2}{\underline{\hspace{1cm}}}} \cdot
            \underset{\text{\small 5th \#}}{\overset{2}{\underline{\hspace{1cm}}}} \cdot
            \underset{\text{\small 6th \#}}{\overset{2}{\underline{\hspace{1cm}}}} \cdot
            \underset{\text{\small 7th \#}}{\overset{2}{\underline{\hspace{1cm}}}} \cdot
            \underset{\text{\small 8th \#}}{\overset{2}{\underline{\hspace{1cm}}}} \cdot
            \underset{\text{\small 9th \#}}{\overset{2}{\underline{\hspace{1cm}}}} \cdot
            \underset{\text{\small 10th \#}}{\overset{2}{\underline{\hspace{1cm}}}}
        \]
        \indent \indent \begin{tabular}[t]{l}
            So, there are $3 \cdot 2^9 = \boxed{\text{1,536}}$ possible ternary sequences.
        \end{tabular} \\
\paragraph{\underline{5.2} \\ \\} 
\textbf{10.} \begin{tabular}[t]{l} 
    What is the probability that an arrangment of the letters $a, b, c, e, f, g$ begins and ends \\
    with a vowel? 
    \end{tabular} \\ \\
    \indent \indent \begin{tabular}[t]{l}
        Since the set of letters $\{a, b, c, e, f, g\}$ has 2 vowels, the first letter has 2 choices and the last \\
        letter has 1 choice. Thus, this gives us P$(2, 2) = 2! = 2$ possible arrangements for the vowels.\\
        The remaining 4 positions have P$(4, 4) $ $= 24$ possible arrangements. The total space of rearr-\\
        anging these six letters is $6! = 720$, making the probability of such an arrangement $\frac{2 \cdot 24}{720} = \boxed{\frac{1}{15}}$
        \end{tabular} \\
    \textbf{32.} \begin{tabular}[t]{l} 
        There are three women and five men who will split up into two four-person teams. How \\
        many ways are there to do this so that there is (at least) one women on each team? 
        \end{tabular} \\ \\
        \indent \indent \begin{tabular}[t]{l}
            First, we want to calculate the total number of ways to put at least one woman on each team. \\
            Since there are 3 women, there are $\binom{3}{1} = 3$ ways to choose at least one woman for a team. \\
            Then, we need to 3 men to fill the remaining spots on the team. There are $\binom{5}{3} = 10$ ways to \\
            choose the remaining men. So, there are $3 \cdot 10 =$ \boxed{\text{30 ways}} to put at least one woman on \\
            each team.
            \end{tabular} \\
    \textbf{40.} \begin{tabular}[t]{l} 
        How many arrangements of 1, 1, 1, 1, 2, 3, 3 are there with the 2 not beside either 3?
        \end{tabular} \\ 
        \indent \indent \begin{tabular}[t]{l}
            The answer to this problem should be equivalent to finding the difference between the total \\
            number of arrangements of the digits and the complement of this problem, the number of \\
            arrangements where 2 is next to at least one 3. Since there are 7 digits with 4 1s, 2 3s and \\
            a 2, 
        \end{tabular}
        \begin{align*}
            \text{P}(7; 4, 1, 2) = \frac{7!}{4! \cdot 1! \cdot 2!} = 105 \text{ ways}
            \end{align*}
        \indent \indent \begin{tabular}[t]{l}
             The complement of the problem has 2 cases:  \\
             \indent \textbf{Case 1}: If 2 is next to at least one 3, the arrangement of 2 and 3 can look like (2,3) \\
             or (3,2). So, there are 2 ways to arrange the 2 and 3. Combining one of the arrangements \\
             into one digit, there'll be 6 digits to arrange: 4 1s, 1 (23) and a 3. So, this case can be arr-\\
             anged in
            \end{tabular} \\
            \begin{align*}
                \text{P}(6; 4, 1, 1) = \frac{6!}{4! \cdot 1! \cdot 1!} = 30 \text{ ways}
                \end{align*}
            \indent \indent \begin{tabular}[t]{l}
                Since the same logic applies for the (3, 2) arrangement, this subcase can be arranged in \\
                $30 \cdot 2 = 60$ ways. \\
                \indent \textbf{Case 2}: If 2 is in between two 3s, the arrangement of the 3s and 2s can look like (3, 2, 3). \\
                Similar to the previous case, combining the (3,2,3) arrangment will lead to arranging 5 digits: \\
                4 1s and 1 (323). So, this case can be arranged in
            \end{tabular} \\
            \begin{align*}
                \text{P}(5; 4, 1) = \frac{5!}{4! \cdot 1!} = 5 \text{ ways}.
            \end{align*}
            \indent \indent \begin{tabular}[t]{l}
                Since the order of the 3s doesn't matter, this is the only arrangment for this case. Since the \\
                case of (3, 2) and (2, 3) is also in (3, 2, 3), the total arrangments for the complement of the \\
                problem is $60 - 5 = 55$ ways. So, the total arrangments with a 2 not beside either 3 is $105 -$\\
                $55$ = \boxed{\text{50 ways}}
                \end{tabular} \\
\end{document}